\scanpage{149}
% archive foliation: 148
Posons
\[
\mu_\alpha=\mu'_\alpha+\mu''_\alpha
\]
$\mu'_\alpha$ partie multiplicative $=$ nbr de branches \uline{formelles} passant par $P_\alpha$, diminué de $1$

$\mu''_\alpha$ partie additive \deleted{\unsure{de}}.
\[
\dim G_{\mathrm{mult}}=\sum\mu'_\alpha-c_0+1=h_1
\]
\margin{\uline{\unsure{invariant topologique}} \uline{des 2 situations}}
[ N.B. Lorsque les $C_i$ \deleted{\unsure{peuvent avoir des}} pts multiples\deleted{\unsure{quelconques}}, alors $h_1$ \unsure{n'est} pas déterminé \deleted{\unsure{par}} combinatoirement par le schém.
\deleted{\unsure{de rencontre}}\inserted{combinatoire de $C$}. Il est en tous cas déterminé par le schém. combinatoire de la \unsure{descente} $\widetilde C\longrightarrow C$; il est aussi le rang du \deleted{$\unsure{\Pi}$} noyau de
$H^1_{\unsure{\ell}}(C)\longrightarrow\deleted{\unsure{\prod}}H^1(\widetilde C)$, (Cohomologie de Weil) ou du conoyau de
\[
\prod\pi_1(\widetilde C_i)-\pi_1(C)\ ]
\]
\[
\dim G_{\mathrm{add}}=\sum\mu''_\alpha
\]
Dire que $\mathrm{Pic}^0$ n'a que des composantes \deleted{\unsure{linéaires}} multiplicatives, signifie que les recollements en $P_\alpha$ sont \unsure{triédraux|triviaux}\inserted{formels; \unsure{Alors} $\mu_\alpha=\mu'_\alpha\geq2$}.
Dire que c'est une V.A. \deleted{signifie} \inserted{implique} que \deleted{\unsure{les}} de plus

